% answers_short.tex —— 2013 年·湖北卷·文科 · 简短答案速查
% 由 tools/make_short_answers.py 自 answers.tex 的 \gksolhead/\gkanswer 提取，
% 并与 2013/2013_数学真题_湖北文_答案要点.md 逐题比对后生成，勿手改（改源 MD 后重跑）。

\begin{gknotice}
\textbf{使用说明}：本节只给最终答案便于快速核对，完整推导见解析版。答案为编写者独立推导，\textbf{非官方标准答案}。
\end{gknotice}

\gkpart{选择题}
\begin{center}\begin{tabular}{ccccccccccc}
\toprule
题号 & 1 & 2 & 3 & 4 & 5 & 6 & 7 & 8 & 9 & 10 \\
\midrule
答案 & B & D & A & D & C & B & A & D & C & B \\
\bottomrule
\end{tabular}\end{center}

\gkpart{填空题}
\begin{itemize}[leftmargin=2.6em]
  \item[\textbf{11.}] $-2+3\mathrm i$。
  \item[\textbf{12.}] (1) $7$；(2) $2$
  \item[\textbf{13.}] $4$。
  \item[\textbf{14.}] $4$。
  \item[\textbf{15.}] $3$。
  \item[\textbf{16.}] $3$。
  \item[\textbf{17.}] (1) $S=3$，$N=1$，$L=6$；(2) $79$
\end{itemize}

\gkpart{解答题}
\begin{itemize}[leftmargin=2.6em]
  \item[\textbf{18.}] (1) $A=\dfrac\pi3$；(2) $\dfrac57$
  \item[\textbf{19.}] (1) $a_n=3\times(-2)^{n-1}$；(2) $\{n\mid n=2k+1,\ k\in\mathbf N,\ k\geqslant5\}=\{11,13,15,\ldots\}$
  \item[\textbf{20.}] (1) 证明见解析；(2) $V_{\text{\text{估}}}<V$
  \item[\textbf{21.}] (1) 当 $a>b$ 时，$f(x)$ 在 $(-\infty,-1)$ 和 $(-1,+\infty)$ 上单调递增；当 $a<b$ 时，$f(x)$ 在 $(-\infty,-1)$ 和 $(-1,+\infty)$ 上单调递减；(2) ① 成等比数列，且 $f\left(\dfrac ba\right)\leqslant f\left(\sqrt{\dfrac ba}\right)$，证明见解析；② 当 $b<a$ 时，$\dfrac ba\leqslant x\leqslant\sqrt{\dfrac ba}$；当 $b>a$ 时，$\sqrt{\dfrac ba}\leqslant x\leqslant\dfrac ba$；当 $a=b$ 时，$x>0$
  \item[\textbf{22.}] (1) $\lambda=1+\sqrt2$；(2) 当 $\lambda>1+\sqrt2$ 时存在；当 $1<\lambda\leqslant1+\sqrt2$ 时不存在
\end{itemize}
